% TOP_PROBLEM: {"id": 30006981, "title": "Breuil-Mezard conjecture: equality of the Galois and automorphic sides", "category_id": 1, "category": {"id": 1, "name": "number_theory", "display_name": "Number Theory", "description": "Properties of integers, prime numbers, Diophantine equations.", "slug": "number-theory", "order_index": 1, "created_at": "2026-07-31T15:26:25.670Z"}, "status": "open"}
% FIRST_POSTED: 2026-09-25
% ENTRY_KIND: statement_only
% !TeX program = lualatex
% Definition and sources only; this file is not an LLM solution attempt.
\documentclass[11pt]{article}
\usepackage[margin=25mm]{geometry}
\usepackage{fontspec}
\setmainfont{Latin Modern Roman}
\usepackage{amsmath,amssymb,mathtools}
\usepackage{unicode-math}
\setmathfont{Latin Modern Math}
\usepackage{xurl,hyperref}
\hypersetup{colorlinks=true,urlcolor=blue}
\setcounter{secnumdepth}{0}
\setlength{\parindent}{0pt}
\setlength{\parskip}{5pt}
\setlength{\emergencystretch}{3em}
\begin{document}
\section*{\texorpdfstring{Breuil-Mezard conjecture: equality of the Galois and automorphic sides}{Breuil-Mezard conjecture: equality of the Galois and automorphic sides}}\label{30006981}
\subsection{Definitions and mathematical statement}
\textbf{Source correction.} The catalogue writes $\sigma_{\mathrm{Gal}}=\sigma_{\mathrm{Aut}}$ as an equality of representations, but its actual source, Volpato's workshop notes, Section 1.1, Conjecture 1, writes a \emph{numerical} equality $\mu_{\mathrm{Gal}}=\mu_{\mathrm{Aut}}$. For a residual two-dimensional representation $\bar\rho$ of $G_{{\mathbb{Q}}_p}$, a Hodge type $k$, and inertial type $\tau$, the Galois quantity is a Hilbert--Samuel multiplicity of a potentially semistable deformation ring modulo a uniformizer. The automorphic expression is
\[
 \mu_{\mathrm{Aut}}(k,\tau,\bar\rho)=\sum_{n,m}a(n,m)\mu_{n,m}(\bar\rho),
\]
where $a(n,m)$ are the multiplicities of weights
$\operatorname{Sym}^n{\mathbb{F}}^2\otimes\det^m$ in the semisimplified reduction of the prescribed type lattice. Hilbert--Samuel multiplicity is the normalized leading coefficient of the local length-growth polynomial. The source's determinant and deformation conditions must be retained. The printed representation-equality target lacks definitions supporting that literal interpretation; its wording is preserved below rather than silently declared equivalent to this numerical conjecture.
\par\smallskip\textit{Formulation and concept references:} \href{https://aimath.org/WWN/padicmodularity/padicmodularity.pdf}{[S1]}.

\subsection{Short English statement}
The cited Breuil–Mézard formulation compares a Galois deformation-ring multiplicity with a weighted automorphic multiplicity, not the undefined representation equality printed in the catalogue.
\subsection{Sources}
\begin{itemize}
\item[S1]AIM workshop problem list, p-adic Representations, Modularity, and Beyond.
\url{https://aimath.org/WWN/padicmodularity/padicmodularity.pdf}.
\textit{Formulation source. Consulted 24 September 2026: Section 1.1, Conjecture 1 is an equality of numerical multiplicities, not the representation equality printed in the catalogue.}
\end{itemize}
\end{document}
